\documentclass[11pt]{article}
\usepackage[margin=1in]{geometry}
\usepackage{amsmath,amssymb}
\usepackage{booktabs}
\usepackage{tabularx}
\usepackage{array}
\usepackage{enumitem}
\usepackage{hyperref}
\usepackage[T1]{fontenc}
\usepackage{caption}
\usepackage{xurl}

\hypersetup{colorlinks=true,linkcolor=blue,urlcolor=blue}
\setlength{\parindent}{0pt}
\setlength{\parskip}{0.55em}
\setlist[itemize]{leftmargin=*,itemsep=0.2em,topsep=0.3em}
\setlist[enumerate]{leftmargin=*,itemsep=0.2em,topsep=0.3em}
\captionsetup{font=small,skip=6pt}
\newcolumntype{Y}{>{\raggedright\arraybackslash}X}
\newcolumntype{P}[1]{>{\raggedright\arraybackslash}p{#1}}

\title{SeQUeNCe 1.3.1 as a host for MUERP\\[0.35em]
{\large \emph{Multi-user Entanglement Routing Design over Quantum Internets}\\
(Zeng et~al., ICDCS 2024)}}
\author{Alexandra Bouillon}
\date{November 10, 2026}

\begin{document}
\maketitle

\begin{abstract}
This note rates \textbf{SeQUeNCe 1.3.1} as a host for MUERP, the routing model in Zeng et~al.\ (ICDCS 2024).
The question is whether SeQUeNCe can reserve exclusive switch memory, follow Algorithms~1--4, and report the one-slot success probability $P$ that Zeng et~al.\ use.

Four measurements, kept separate:
(i)~\textbf{admission} on Zeng et~al.\ Fig.~4a (one switch in the middle, three users around it);
(ii)~Algorithms~2--4 trees written into SeQUeNCe forwarding tables on one Waxman graph with 10 users and 50 switches, plus two groups sharing switch memory $Q_r$;
(iii)~entanglement-generation and swapping fidelity on that same Fig.~4a layout; Eq.~(2) computed in \texttt{muerp.eq2}, which is not a SeQUeNCe metric;
(iv)~a linear repeater chain, used only to time the simulator.

Must-have scores $[3,5,5,4,4,3]$; weighted mean $4.00$. No Must-have is 1.
\end{abstract}

\tableofcontents
\vspace{1em}

%----------------------------------------------------------------------
\section{What is being scored}

The object under test is \textbf{SeQUeNCe 1.3.1}.
The task is: can it host MUERP as stated by Zeng et~al.~\cite{zeng2024}?

\subsection{Yardstick: MUERP (Zeng et~al.)}

Zeng, Zhang, Shang, Liu, Liu, and Yang~\cite{zeng2024} define multi-user entanglement routing (MUERP):
users $U$ and switches $R$ on $G=(V,E)$ with $V=U\cup R$; an entanglement tree of $|U|-1$ pairwise channels; per-switch budget $Q_r$; two exclusive qubits per channel through a switch.
Fiber is treated as abundant; switches are the bottleneck.

Per-channel success in one synchronized slot is
\begin{equation}
P_A = q^{\ell-1}\prod_{i} p_{i,i+1},\qquad
p_{i,i+1}=\exp(-\alpha L_{i,i+1}),
\label{eq:pa}
\end{equation}
with uniform swap success $q$ and default $\alpha=10^{-4}$ (length unit $\approx 1$\,km).
The tree rate is
\begin{equation}
P=\prod_{A} P_A,
\label{eq:p}
\end{equation}
and $P=0$ if any required channel cannot be reserved.
Routing is offline: Algorithms~1--4 choose the tree, then every link and swap is attempted in the same slot.

Zeng et~al.\ Sec.~V-A evaluates this model on random graphs (Waxman, Watts--Strogatz, Volchenkov) with $|U|=10$, $|R|=50$, $D=6$, $Q_r=4$, $q=0.9$, area $10^4\times 10^4$, and 20 replicas.
They also name two extensions measured here: fidelity decay, and concurrent groups sharing one $Q_r$ pool.
Algorithms~1--4 are their routing contribution.
SeQUeNCe is scored on whether it can hold that model.

\paragraph{Four-node layout versus entanglement tree.}
Zeng et~al.\ Fig.~4a is a \textbf{four-node layout}: one switch $S$ in the middle, and three users (Alice, Bob, Carol) around it.
That is the device graph.
The \textbf{entanglement tree} on that graph is the two channels Alice--$S$--Bob and Alice--$S$--Carol.
This note uses \emph{four-node layout} for the device graph and \emph{tree} for the channels.

\subsection{What SeQUeNCe is}

SeQUeNCe~\cite{sequence} 1.3.1 (\texttt{pip install sequence}, isolated \texttt{SeQuEnCe/.venv}, Python 3.12.4, \texttt{import sequence as sq}) is a discrete-event simulator with a resource-reservation protocol (RSVP), static routing, and \texttt{RequestApp}.

Two numbers in this note must not be mixed:
\begin{itemize}
\item \textbf{Slot $P$} --- Eqs.~\eqref{eq:pa}--\eqref{eq:p}, one simultaneous attempt of every link and swap.
\item \textbf{SeQUeNCe throughput} --- \texttt{RequestApp.memory\_counter} over a service window, with retries.
\end{itemize}
Zeng et~al.\ report slot $P$.
SeQUeNCe attempts are sequential in simulated time.
The slot number reported here is RSVP admission followed by Eq.~\eqref{eq:p} in \texttt{muerp.eq2}.

\subsection{Measurements}

\begin{center}
\begin{tabularx}{\textwidth}{@{}P{1.55in} Y Y@{}}
\toprule
Measurement & What ran & Role \\
\midrule
Four-node admission & RSVP on Fig.~4a; analytic $P$ & Exclusive $Q_r$ (M2) \\
Four-node physics & Entanglement generation and swapping on Fig.~4a & Throughput vs $P$; fidelity (M11) \\
Four-node Alg.~4 paths & Alg.~4 next hops in forwarding tables; cutoff on/off & Inject routing (M5), timing (M10) \\
Waxman load & Dummy spanning tree, 10 users + 50 switches & Size at Zeng et~al.'s node count \\
Waxman Alg.~2--4 paths & Those algorithms' trees on the same scale; two groups sharing $Q_r$ & M5 at 10+50; concurrent groups (M12) \\
Linear chain & 2--100 routers, $0.2$\,s generation window & Simulator wall-clock (M14) \\
\bottomrule
\end{tabularx}
\end{center}

Limits of these runs:
Algorithms~1--4 choose paths; RSVP chooses which memories to reserve;
attempts are sequential, not one simultaneous slot;
the linear chain times the simulator, it is not a MUERP instance.

%----------------------------------------------------------------------
\section{Rating method}

Scores 1--5: 5 native, 4 modest glue, 3 reimplement a primitive, 2 breaks an assumption, 1 not practical.
Must (M1--M6) $w=3$; important (M7--M10) $w=2$; nice (M11--M15) $w=1$.
Any Must-have of 1 bars SeQUeNCe as a MUERP host.
Composite $\sum w_i s_i / \sum w_i$.

The gate is Fig.~4a.
This note uses $p=q=0.9$.
Case A: $Q_S=4$ $\Rightarrow$ $P=(p^2 q)^2=0.531441$.
Case B: $Q_S=2$ $\Rightarrow$ $P=0$.
Failing either case fails M2/M6.

%----------------------------------------------------------------------
\section{Installation}

Isolated venv; \texttt{sequence} 1.3.1; timeline smoke test exited 0.
Python is x86\_64 under Rosetta on arm64 macOS.

Physics runs used
\begin{verbatim}
sq.EntanglementGenerationA.set_global_type(SINGLE_HERALDED)
sq.EntanglementGenerationB.set_global_type(SINGLE_HERALDED)
sq.EntanglementSwappingA.set_formalism(
    sq.constants.BELL_DIAGONAL_STATE_FORMALISM)
sq.EntanglementSwappingB.set_formalism(
    sq.constants.BELL_DIAGONAL_STATE_FORMALISM)
\end{verbatim}
Swapping formalism must match the Bell-diagonal manager.

%----------------------------------------------------------------------
\section{Algorithms 1--4}

Implemented in \texttt{muerp/algorithms.py} from Zeng et~al.'s descriptions.
They run in Python, outside SeQUeNCe, and return trees whose next hops are then written into SeQUeNCe forwarding tables (Sec.~\ref{sec:inject}, Sec.~\ref{sec:kernel60}).

\begin{itemize}
\item Alg.~1: Dijkstra on $\alpha L-\ln q$; skip remaining $Q_r<2$; $\mathrm{RATE}=q^{\ell-1}\prod p$.
\item Alg.~2: all-pairs Alg.~1 with infinite $Q_r$, then Kruskal. Does not enforce finite $Q_r$ on the chosen tree.
\item Alg.~3: greedy insert under remaining $Q_r$, then reconnect.
\item Alg.~4: Prim growth with Alg.~1.
\end{itemize}

%----------------------------------------------------------------------
\section{Four-node layout: admission}
\label{sec:fig4a}

Zeng et~al.\ Fig.~4a: one switch in the middle, three users around it.
In this run only: \texttt{router\_0}${=}S$, \texttt{router\_1}${=}$Alice, \texttt{router\_2}${=}$Bob, \texttt{router\_3}${=}$Carol.
On the Waxman runs, \texttt{router\_0}--\texttt{router\_9} are users (same names, different runs).
Admission only: timeline stop $0.1$\,s, before $t=1$\,s.

\begin{table}[h]
\centering
\caption{Four-node layout admission in SeQUeNCe 1.3.1. Analytic $P$, not throughput.}
\label{tab:fig4a}
\begin{tabular}{@{}lcccc@{}}
\toprule
Case & $Q_S$ & Admitted & Memories on $S$ & $P$ \\
\midrule
A & 4 & both channels & $4/4$ & $0.531441$ \\
B & 2 & Alice--Bob only; Alice--Carol rejected & $2/2$ & $0$ \\
\bottomrule
\end{tabular}
\end{table}

RSVP takes two exclusive memories per channel and refuses the second when $Q_S=2$.

Do not read one \texttt{RequestApp.reservation\_result} for two starts on Alice; use \texttt{accepted\_reservations} on the switch.

%----------------------------------------------------------------------
\section{Four-node layout: device physics versus slot $P$}
\label{sec:physics}

Bernoulli Monte Carlo, $2\times 10^4$ trials, same two channels: $0.534\pm 0.004$ (s.e.) vs analytic $0.531441$.
That only checks the product formula.

SeQUeNCe entanglement generation, $Q_S=4$, $0.2$\,s window after $t=1$\,s: Alice--Bob $7$ pairs, Carol--Alice $10$, wall $0.07$\,s.
That is throughput, not $P$.

Attenuation was set so each 5\,km half-link has transmittance $0.9$, while the NetworkX graph used for delay is $10$\,m.
Admission is unaffected.
Generation is a device stack calibrated to $p$, not a literal $\exp(-\alpha L)$ slot.

\paragraph{Fidelity events ($T=0.05$\,s).}
Entanglement-generation fidelity mean $0.90$ ($n=72$).
Swapping fidelity mean $0.57$, min $0.39$ ($n=16$).
That is a swap-fidelity log, not a per-channel trajectory object.
\texttt{swapping\_degradation} is not a keyword on Bell-diagonal swapping in 1.3.1.

%----------------------------------------------------------------------
\section{Four-node layout: Alg.~4 paths, then cutoff}
\label{sec:inject}

In Python, Alg.~4 builds the tree.
SeQUeNCe forwarding tables are cleared and set to those next hops (and the reverse).
One \texttt{RequestApp} per channel; the second initiator is Carol.

\begin{table}[h]
\centering
\caption{Alg.~4 paths in RSVP on the four-node layout. Alg.~4's product at $L=10$\,m is $0.807$; Fig.~4a slot $P$ at $p=0.9$ is $0.531441$.}
\begin{tabular}{@{}lcccc@{}}
\toprule
$Q_S$ & Alg.~4 tree & RSVP & Paths match tables? & Slot $P$ at $p=0.9$ \\
\midrule
4 & $1$--$0$--$2$ and $1$--$0$--$3$ & both accepted & yes & $0.531441$ \\
2 & empty & no starts & n/a & $0$ \\
\bottomrule
\end{tabular}
\end{table}

RSVP chooses memories, so M5 is 4 rather than 5.
Alg.~4's own product on the NetworkX graph used $L=10$\,m, so $p\approx 1$ and that product is $0.807$ (\texttt{alg4\_P} in \texttt{inject\_alg4\_coherence.json}).
Use $0.531$ for the $p=0.9$ instance.

\begin{table}[h]
\centering
\caption{Physics on the Alg.~4 tree, $Q_S=4$, $0.2$\,s window. Throughput, not slot $P$.}
\label{tab:coh}
\begin{tabular}{@{}lcccc@{}}
\toprule
Setting & Gen.\ / swap success & Expired & Alice, Carol counters & Wall \\
\midrule
\texttt{coherence\_time}${=}{-}1$, cutoff off & 80 / 18 & 0 & 7, 10 & $0.068$\,s \\
$T=0.01$\,s, cutoff on & 84 / 13 & 44 & 0, 0 & $0.067$\,s \\
\bottomrule
\end{tabular}
\end{table}

Bob's counter is 0 (responder).
Admission is unchanged when cutoff is on; retry-throughput changes.
A timed simulator can disable wait (infinite lifetime) or kill pairs (short cutoff).
That is not Zeng et~al.'s simultaneous slot.

%----------------------------------------------------------------------
\section{Waxman, 10 users and 50 switches}

Two runs at Zeng et~al.'s node count.
They use different trees; do not mix the two $P$ values.
On both runs, \texttt{router\_0}--\texttt{router\_9} are users.

\subsection{Load: dummy spanning tree}
\label{sec:waxman}

Waxman, 60 nodes, area $10^4\times 10^4$, seed 0, NetworkX $\alpha=0.15$, $\beta=0.8$, mean degree $6.33$ (Zeng et~al.\ use $D=6$).
$Q_r=4$ on switches.
Tree: a minimum spanning tree of shortest paths among users (9 channels).
This tree is a load generator, not Alg.~2, 3, or 4.

Admitted $9/9$; analytic $P=0.0352$; max occupancy $2/4$; wall $0.55$\,s (of which $0.50$\,s is \texttt{RouterNetTopo}); peak RSS $160$\,MB.
Admission only: no entanglement generation.

\subsection{Algorithms 2--4 paths}
\label{sec:kernel60}

\texttt{muerp/showcase\_kernel.py}, seed 0, $|U|=10$, $|R|=50$, $Q_r=4$.
\texttt{muerp/algorithms.py} builds the tree; forwarding tables are set to those next hops; RSVP admits; \texttt{muerp.eq2.tree\_rate} computes $P$.
SeQUeNCe has no native Eq.~(2).

\begin{table}[h]
\centering
\caption{One Waxman 10+50, admission only. Algs.~2--4 coincided on this draw. $P=0.0197$ is this tree, not the dummy tree in Sec.~\ref{sec:waxman}.}
\label{tab:kernel60}
\begin{tabular}{@{}lcccc@{}}
\toprule
Solver & Alg.~$P$ & Admitted & Paths match tables? & RSVP Eq.~(2) \\
\midrule
Alg.~2, 3, 4 & 0.0197 & $9/9$ & yes & 0.0197 \\
\bottomrule
\end{tabular}
\end{table}

Wall $\approx 0.5$\,s per solver (mostly \texttt{RouterNetTopo}). Max occupancy $2/4$.
One replica; admission only.

\paragraph{Two groups sharing $Q_r$.}
Users $\{0..4\}$ and $\{5..9\}$, Alg.~4 each, one RSVP pool.
At $Q_r=4$ and $Q_r=2$ both groups admitted $4/4$ on this draw (max occupancy 2; at $Q_r=2$, six switches saturated).
The pool is shared. These two trees did not block each other.

%----------------------------------------------------------------------
\section{Linear repeater chain}
\label{sec:chain}

\texttt{probes/repeater\_chain\_speed.py}: path of $n$ \texttt{QuantumRouter}s, one \texttt{RequestApp} from \texttt{router\_0} to the far end, $q=0.9$, 8 memories, infinite coherence.
Admission stops before the service window.
Physics uses the same $0.2$\,s window as the four-node layout.
Each size ran in its own process.

This is SeQUeNCe wall-clock on a line, not a MUERP random graph and not slot $P$.

\begin{table}[h]
\centering
\caption{Linear chain, SeQUeNCe 1.3.1. RSVP admitted every size.}
\label{tab:chain}
\begin{tabular}{@{}rrrrrr@{}}
\toprule
Routers & Hops & Admit (s) & Physics (s) & Peak RSS (MB) & Pairs in $0.2$\,s \\
\midrule
2 & 1 & 0.005 & 0.019 & 155 & 1 \\
5 & 4 & 0.007 & 0.036 & 156 & 1 \\
10 & 9 & 0.011 & 0.066 & 157 & 0 \\
20 & 19 & 0.027 & 0.155 & 155 & 0 \\
50 & 49 & 0.18 & 0.60 & 159 & 0 \\
100 & 99 & 1.02 & 2.02 & 167 & 0 \\
\bottomrule
\end{tabular}
\end{table}

At 50--100 routers, most of the wall-clock is \texttt{RouterNetTopo} construction, not the $0.2$\,s of entanglement generation.
End-to-end \texttt{memory\_counter} is already 0 at 10 hops in this window: swap-wait along the chain.

%----------------------------------------------------------------------
\section{API map (M1--M6)}

Hooks are from installed 1.3.1 (\texttt{set\_swapping\_success\_rate} is absent).

\paragraph{M1 (score 3).}
\texttt{generate\_config} $\to$ \texttt{RouterNetTopo}.
Only \texttt{QuantumRouter} and \texttt{BSMNode}.
\texttt{node\_type} is dropped.
No User class; \texttt{sq.Switch} is optics.
Tag roles by name and \texttt{memo\_size}.

\paragraph{M2 (score 5).}
Intermediates reserve $2\times$\texttt{memory\_size}.
\texttt{MemoryTimeCard} refuses overlap.
Fig.~4a post-edits $S$ to 4 or 2.

\paragraph{M3 (score 5).}
\texttt{swapping\_success\_prob} (default 1) is Zeng et~al.'s $q$, not mid-link \texttt{SingleHeraldedBSM}.

\paragraph{M4 (score 4).}
Loss is dB/m, not $\exp(-\alpha L)$.
Convert, or multiply $p$ after admission.

\paragraph{M5 (score 4).}
\texttt{RequestApp.start} has no path argument.
Forwarding tables can be overwritten.
RSVP chooses memories.
Evidence: Algs.~2--4 paths on one 10+50 Waxman (Sec.~\ref{sec:kernel60}).

\paragraph{M6 (score 3).}
No native $P$.
Reject $\Rightarrow$ $P=0$; else multiply $p$ and $q$ in the script.
Do not use \texttt{memory\_counter}.

%----------------------------------------------------------------------
\section{Scorecard}

M10 is an interpolation: anchors are 3 (timed only) and 5 (slotted \emph{and} timed).
M14 is SeQUeNCe wall-clock on the graphs in this note.

\begin{table}[h]
\centering
\caption{SeQUeNCe as a MUERP host.}
\label{tab:score}
{\small
\begin{tabularx}{\textwidth}{@{}l P{1.2in} r r Y@{}}
\toprule
ID & Metric & $w$ & $s$ & Evidence / limit \\
\midrule
M1 & Users vs switches & 3 & 3 & Same \texttt{QuantumRouter} type. \\
M2 & Exclusive $Q_r$ & 3 & 5 & Fig.~4a: $Q_S=4$ admits, $Q_S=2$ rejects. \\
M3 & BSM, constant $q$ & 3 & 5 & \texttt{swapping\_success\_prob}. \\
M4 & $p=\exp(-\alpha L)$ & 3 & 4 & dB/m; convert or compute $p$ outside. \\
M5 & Inject Algs~1--4 & 3 & 4 & Algs.~2--4 paths on one 10+50 Waxman; RSVP chooses memories. \\
M6 & Export $P$ & 3 & 3 & \texttt{muerp.eq2}. No SeQUeNCe metric named $P$. \\
M7 & Independent channels & 2 & 4 & Two channels, 4 exclusive qubits on $S$. \\
M8 & Multi-user tree & 2 & 3 & Pairwise apps; product outside. \\
M9 & Scale & 2 & 4 & 10+50 RSVP $\approx 0.5$\,s, one replica. \\
M10 & Slot vs decoherence & 2 & 4 & Timed; cutoff on/off. Slot $P$ is Eq.~(2) after RSVP. \\
M11 & Fidelity extension & 1 & 4 & Swapping fidelity mean $0.57$ vs generation $0.90$ on the four-node layout. \\
M12 & Concurrent groups & 1 & 4 & Two 5-user trees, one $Q_r$ pool; both fit on this draw. \\
M13 & Open internals & 1 & 5 & PyPI 1.3.1 modules. \\
M14 & Runtime on this workload & 1 & 4 & Waxman $0.55$\,s / $160$\,MB; chain $\le 2$\,s at 100 routers. \\
M15 & State representation & 1 & 5 & \texttt{bell\_diagonal} for this stack. \\
\midrule
 & Mean & 31 & 4.00 & Must-have scores $[3,5,5,4,4,3]$. \\
\bottomrule
\end{tabularx}}
\end{table}

Numerator $3(3+5+5+4+4+3)+2(4+3+4+4)+(4+4+5+4+5)=124$. Mean $124/31=4$.

%----------------------------------------------------------------------
\section{Implementation pitfalls}

\begin{enumerate}
\item \texttt{reservation\_result} is overwritten by a second start on the same app.
\item Reverse the second channel if both would share Alice as initiator.
\item Swapping formalism must match the Bell-diagonal manager.
\item \texttt{generate\_config} ignores \texttt{node\_type} and writes one \texttt{memo\_size}.
\item Default \texttt{swapping\_success\_prob} is 1, not $q=0.9$.
\item Mid-link \texttt{BSMNode} is not Zeng et~al.'s switch and not $q$.
\item Slot $P$ $\neq$ \texttt{memory\_counter}. Stop before the service window for admission.
\item \texttt{alg4\_P}${}=0.807$ in the injection JSON is the $10$\,m product, not $0.531$.
\item \texttt{router\_0} is $S$ on the four-node layout and a user on the Waxman runs.
\item Matplotlib font cache can look like a hang. Set \texttt{MPLCONFIGDIR=/tmp/sequence-mpl}.
\end{enumerate}

%----------------------------------------------------------------------
\section{Files}

Paths relative to \texttt{MUERP/Simulators/}.

\begin{center}
\begin{tabular}{@{}ll@{}}
\toprule
File & Role \\
\midrule
\texttt{sequence\_api\_map\_m1\_m6.md} & Calls, no scores \\
\texttt{probes/fig4a\_sequence.py} & Four-node layout admission \\
\texttt{probes/waxman\_replica\_sequence.py} & Dummy-tree 10+50 RSVP \\
\texttt{probes/repeater\_chain\_speed.py} & Linear chain 2--100 routers \\
\texttt{muerp/results/repeater\_chain\_speed.json} & Chain wall-clock dump \\
\texttt{muerp/algorithms.py} & Algorithms~1--4 \\
\texttt{muerp/eq2.py} & Eq.~(1)--(2) \\
\texttt{muerp/showcase\_kernel.py} & Algs~2--4 paths; two groups; four-node fidelity \\
\texttt{muerp/mc\_physics.py} & Slot Monte Carlo + unpaired generation \\
\texttt{muerp/inject\_alg4\_coherence.py} & Alg.~4 paths; cutoff \\
\texttt{SeQuEnCe/.venv} & 1.3.1 \\
\bottomrule
\end{tabular}
\end{center}

\begin{verbatim}
export MPLCONFIGDIR=/tmp/sequence-mpl
export PYTHONPATH=MUERP/Simulators
SeQuEnCe/.venv/bin/python -m muerp.inject_alg4_coherence
\end{verbatim}

%----------------------------------------------------------------------
\section{What these runs show}

SeQUeNCe 1.3.1 can reserve exclusive qubits the way MUERP requires on Zeng et~al.'s four-node layout, refuse a second channel when $Q_S=2$, and follow Alg.~2/3/4 next-hop tables on a 10+50 Waxman ($9/9$ admitted, $P=0.0197$ on that tree).
Eqs.~\eqref{eq:pa}--\eqref{eq:p} are computed in \texttt{muerp.eq2}, not by SeQUeNCe.
Two 5-user groups can share one $Q_r$ pool; on this draw both completed.
Swapping events carry fidelity (mean $0.57$ after swaps vs $0.90$ at generation).
Attempts in SeQUeNCe are sequential; the simultaneous slot is Eq.~\eqref{eq:p} after RSVP.

Must-have scores $[3,5,5,4,4,3]$, weighted mean $4.00$.
SeQUeNCe can host the MUERP admission model on the instances in this note.

\begin{thebibliography}{9}
\bibitem{zeng2024}
Y.\ Zeng, J.\ Zhang, X.\ Shang, J.\ Liu, Z.\ Liu, and Y.\ Yang,
``Multi-user Entanglement Routing Design over Quantum Internets,''
in \textit{Proc.\ IEEE 44th International Conference on Distributed Computing Systems (ICDCS)}, 2024.
DOI: \href{https://doi.org/10.1109/ICDCS60910.2024.00033}{10.1109/ICDCS60910.2024.00033}.
Repo source: \texttt{MUERP/MultiUser\_Paper\_1.md}.

\bibitem{sequence}
X.\ Wu et~al.,
``SeQUeNCe: a customizable discrete-event simulator of quantum networks,''
\textit{Quantum Sci.\ Technol.}, 2021.
Package: \texttt{sequence} 1.3.1, \url{https://github.com/sequence-toolbox/SeQUeNCe}.
\end{thebibliography}

\end{document}
